\begin{proof}
La afirmación que hay que demostrar para obtener (7.4.4.1) se deduce
de este lema observando que, puesto que $B$ es íntegro, el punto
genérico de $X$ pertenece a $U$.

Se observará que, como $C$ es una $A$-álgebra de tipo finito y $\mathfrak n$
un ideal maximal de $C$, el cuerpo residual $k'$ de
$C_{\mathfrak n}=B$ (y, por tanto, también de $B'$) es una extensión
\emph{finita} de $k$ (\textbf I, 6.4.11 y 6.4.2).

Si se pone $Y=\operatorname{Spec}(C)$, de (6.12.8) resulta que
$\operatorname{Reg}(Y)$ es \emph{abierto} en $Y$; como los anillos locales de
$X$ son anillos locales de $Y$ (\textbf I, 2.4.2), se tiene
$U=X\cap\operatorname{Reg}(Y)$, y por tanto $U$ es \emph{abierto} en $X$; por
otra parte, según (6.12.7), $U'$ es abierto en $X'$, de modo que
$S'=X'\setmin U'$ es cerrado. En consecuencia, $S'\cap f^{-1}(U)$ es
\emph{localmente cerrado} en $X'$, y basta demostrar que este conjunto es
\emph{vacío}. Se sabe por (5.1.10) que, en caso contrario, existiría en
$S'\cap f^{-1}(U)$ un ideal primo $\mathfrak p'$ tal que
$\dim(B'/\mathfrak p')\leqslant1$. Observemos primero que
$\mathfrak p'$ no puede ser el ideal maximal $\mathfrak mB'$ de $B'$,
donde $\mathfrak m$ es el ideal maximal de $B$. En efecto, esto
significaría que $B=B_{\mathfrak m}$ sería regular, y por tanto también
$B'=\widehat B$ (\textbf{0}, 17.1.5); se tendría entonces $\mathfrak mB'\in U'$,
en contra de la hipótesis. Debe ser, pues,
$\dim(B'/\mathfrak p')=1$. Pongamos
$\mathfrak p=B\cap\mathfrak p'$; por hipótesis, $B_{\mathfrak p}$ es
regular, pero $B'_{\mathfrak p'}$ no lo es; como
$B'_{\mathfrak p'}$ es un $B_{\mathfrak p}$-módulo plano, se deduce de
(6.5.2) que la fibra $Z$ de $f$ en el punto $\mathfrak p$ \emph{no es regular
en el punto} $\mathfrak p'$. Mostremos que \emph{es posible reducirse al caso en
que} $\mathfrak p=0$. En efecto, en el caso general, si se pone
$\mathfrak q=\mathfrak p\cap C$, $\mathfrak r=\mathfrak p\cap A$,
$C/\mathfrak q$ es una $(A/\mathfrak r)[\bar t]$-álgebra finita (donde
$\bar t$ es la clase de $t$ módulo $\mathfrak q$); como
$\mathfrak p=\mathfrak qB$, $B/\mathfrak p$ es igual a
$(C/\mathfrak q)_{\mathfrak n/\mathfrak q}$, y
$\mathfrak n/\mathfrak q$ es un ideal maximal de $C/\mathfrak q$
situado sobre el ideal maximal $\mathfrak m/\mathfrak r$ de
$A/\mathfrak r$. Se ve así que $A/\mathfrak r$, $C/\mathfrak q$ y
$B/\mathfrak p$ siguen satisfaciendo las hipótesis de (7.4.4.2); como
el completado de $B/\mathfrak p$ es $B'/\mathfrak pB'$, esto demuestra
nuestra afirmación. Supongamos, pues, $\mathfrak p=0$, de manera que
$Z$ es la fibra genérica y el homomorfismo $B\to B'/\mathfrak p'$ es
\emph{inyectivo}. Pongamos $V=B'/\mathfrak p'$ y distingamos dos casos:

I) \emph{$V$ es una $A$-álgebra finita.} --- Como $B\subset V$, $B$ es
\emph{a fortiori} una $A$-álgebra finita; y, puesto que $A$ es completo,
lo mismo ocurre con $B$ (Bourbaki, \emph{Alg. comm.}, cap. IV, \S~2,
n\textsuperscript{o}~5, cor. 3 de la prop. 9). De ahí que $B'=B$ y
$\mathfrak p'=0$, por lo que $B'_{\mathfrak p'}$ es un cuerpo y, por
consiguiente, un anillo regular, en contra de la hipótesis.

II) \emph{$V$ no es una $A$-álgebra finita.} --- Como el anillo local
$A$ es \emph{completo}, esto implica que $V$ no es una $A$-álgebra
\emph{cuasifinita} ($\mathbf{0}_{\mathrm I}$, 7.4.1 y 7.4.2); pero, por hipótesis, el
cuerpo residual $k'$ de $V$ es una extensión \emph{finita} del cuerpo residual
$k$ de $A$,
\oldpage[IV-2]{202}
por tanto, según ($\mathbf{0}_{\mathrm I}$, 7.4.4), el ideal $\mathfrak mV$
\emph{no es un ideal de definición de} $V$. Como $V$ es un anillo local noetheriano
íntegro de dimensión 1, 0 es el único ideal de $V$ que no es un ideal
de definición; por ello, $\mathfrak mV=0$. Pero $A\subset V$ y $V$ es
íntegro, de donde $\mathfrak m=0$ y $A=k$ es un \emph{cuerpo}. Se deduce ante
todo que $\dim(C)\leqslant1$ (\textbf 0, 16.1.5); pero, como $\dim(V)=1$, las
relaciones
$\dim(V)\leqslant\dim(B')=\dim(B)\leqslant\dim(C)$
muestran que esto implica
$\dim(C)=\dim(B)=\dim(B')=\dim(B'/\mathfrak p')=1$ y, en consecuencia,
$\mathfrak p'$ es necesariamente un \emph{ideal primo minimal} de $B'$.
Llegaremos, pues, también a una contradicción si demostramos que
$B'_{\mathfrak p'}$ es un \emph{cuerpo}, o, lo que es lo mismo, que el anillo
$B'$ es reducido. Ahora bien, como $C$ es una $k$-álgebra de tipo finito,
la clausura integral $C_1$ de $C$ es una $C$-álgebra \emph{finita} (Bourbaki,
\emph{Alg. comm.}, cap. V, \S~3, n\textsuperscript{o}~2, teorema 2);
si se pone $S=C\setmin\mathfrak n$, $B_1=S^{-1}C_1$ es la clausura integral de
$B$, y por tanto $B_1$ es una $B$-álgebra \emph{finita}, y en consecuencia es
un anillo semilocal noetheriano, íntegro, integralmente cerrado y de
dimensión 1 (\textbf{0}, 16.1.5). Si $\mathfrak m_j$ $(1\leqslant j\leqslant h)$
son sus ideales maximales, los $(B_1)_{\mathfrak m_j}$ son, por tanto,
anillos de valoración discreta (\textbf{II}, 7.1.6), y el completado $B'_1$ de
$B_1$ es el producto directo de los anillos de valoración discreta
completados de los $(B_1)_{\mathfrak m_j}$ (Bourbaki, \emph{Alg. comm.},
cap. III, \S~2, n\textsuperscript{o}~13, prop. 18); $B'_1$ es, pues,
reducido, y como el completado $B'$ de $B$ es un subanillo de $B'_1$
(Bourbaki, \emph{Alg. comm.}, cap. IV, \S~2, n\textsuperscript{o}~5,
cor. 3 de la prop. 9), también es un anillo reducido. Queda demostrado.
\end{proof}
